\documentclass[10pt,a4paper]{amsart}
\usepackage[margin=19mm]{geometry}
\usepackage{amsmath,amssymb,mathtools}
\usepackage[scr=rsfs,cal=euler]{mathalfa}
\usepackage{xfrac}
\usepackage[unicode,hidelinks,bookmarks=false]{hyperref}
\newcommand{\ie}{i.e.\ }
\newcommand{\cf}{cf.\ }
\newcommand{\resp}{respectively\ }
\newcommand{\Subsection}{\subsection}
\providecommand{\nobreakdash}{\nobreak\hbox{-}}
\setlength{\emergencystretch}{2em}
\multlinegap0pt
\usepackage{bm}
\usepackage{bbm}
\usepackage{dsfont}
\usepackage{xspace}
\usepackage{cancel}
\usepackage{mathtools}
\usepackage{amsthm}
\makeatletter
\@ifundefined{theorem}{\newtheorem{theorem}{Theorem}}{}
\@ifundefined{lemma}{\newtheorem{lemma}[theorem]{Lemma}}{}
\makeatother
\newcommand{\N}{\mathbb{N}}
\newcommand{\colonequal}{\mathrel{\mathop:}=}
\title[Universal Matrices for Counting Fibonomial and $C$-nomial Co]{Universal Matrices for Counting Fibonomial and $C$-nomial Coefficients by their $p$-adic Valuations}
\author{Arav Chand}
\date{Source: arXiv:2508.18461v1 (CC BY 4.0)}
\begin{document}
\maketitle

\noindent\emph{Source and licence.} Universal Matrices for Counting Fibonomial and $C$-nomial Coefficients by their $p$-adic Valuations. Version arXiv:2508.18461v1 (CC BY 4.0), distributed under the Creative Commons Attribution 4.0 International licence, as recorded in the arXiv licence metadata for this exact version and shown on the abstract page https://arxiv.org/abs/2508.18461v1. Full rights notice: licenses/arxiv-2508.18461-CC-BY-4.0.txt.\par\smallskip

\noindent\textit{Editorial scope.} Counted unit: the Lemma whose source label is \texttt{lemma:idealCnomialSummationInBinomials} in the article (source0.tex lines 441--445, source label \texttt{lemma:idealCnomialSummationInBinomials}), delivered together with the article's own complete original proof of it (source0.tex lines 447--488, one \texttt{proof} environment of 42 lines). No mathematics is written, repaired or re-derived: statement and proof are the article's own text line for line. Required context retained uncounted: lemma:idealCorial (lines 412--418). Every definition, display and label of those lines is retained unchanged. Omitted from this package and not redistributed: 1--411, 419--440, 446--446, 489--692. Nothing inside a retained excerpt is omitted or altered: every retained body is delivered exactly as the article's lines are written, so no omission receipt of any kind accompanies this package. Citations: the retained excerpts carry no citation command at all, so no bibliography entry of the article is transcribed and no reference is redistributed. Reference adaptation: none was needed and none was made; every reference inside the delivered unit resolves inside it, so no unresolved reference or citation is produced. Printed slips kept literal: no printed word, symbol, hypothesis or number of the counted statement or of its proof was altered. Layout and class: the article's own class is \texttt{amsart} with packages including amsmath, amssymb, graphicx, latexsym, color, stmaryrd, hyperref; the wrapper is the lane's pinned \texttt{amsart} editorial wrapper at 10pt on A4, which restates the theorem-like environments the retained text needs and adds the article's own macro definitions that the retained text needs (\texttt{\textbackslash N}, \texttt{\textbackslash colonequal}), with extra wrapper packages (bm, bbm, dsfont, xspace, cancel, mathtools). The delivered numbering is the unit's own. Assets: no retained span contains a figure environment, a graphics-inclusion command, a PDF-page inclusion, a table or a TikZ picture; no image file is copied into this package, the package declares no asset dependency and no image file is redistributed with it. Licence: version arXiv:2508.18461v1 of this article is distributed under the Creative Commons Attribution 4.0 International licence, as recorded in the arXiv licence metadata for this exact version and shown on the abstract page https://arxiv.org/abs/2508.18461v1. A full rights notice is delivered at licenses/arxiv-2508.18461-CC-BY-4.0.txt. Characters: every retained excerpt is pure ASCII, so the delivered engine, XeLaTeX with its default Latin Modern text font, reports no missing character.
\begin{lemma} % [$C$-orial Legendre's Formula for Ideal Primes]
\label{lemma:idealCorial}
Let $ p $ be an ideal prime for a strong divisibility sequence $ C $. Then for any integer $ n \geq 0 $,
\[
\nu_p(n!_C) = s(p) \left\lfloor \frac{n}{\alpha(p)} \right\rfloor + \nu_p\!\left( \left\lfloor \frac{n}{\alpha(p)} \right\rfloor ! \right).
\]
\end{lemma}

\begin{lemma}
\label{lemma:idealCnomialSummationInBinomials}
Let $p$ be an ideal prime for a strong divisibility sequence $C$. For any integer $n\geq0$, let $n=\alpha(p)n'+r$ where $0 \leq r <\alpha(p)$ and $ n' \in \N $. Then
$$\sum_{m=0}^{n} x^{\nu_p\!\left( \binom{n}{m}_C \right)} = (r+1)\sum_{m'=0}^{n'} x^{\nu_p\!\left(\binom{n'}{m'}\right)} + (\alpha(p) - r - 1)x^{s(p) - 1}\sum_{m'=0}^{n'-1} x^{1+\nu_p\!\left(n'\binom{n'-1}{m'}\right)}.$$
\end{lemma}

\begin{proof}
We first analyze the exponent of the left-hand-side $$\nu_p\!\left(\binom{n}{m}_C\right) = \nu_p(n!_C)-\nu_p(m!_C)-\nu_p\!\left((n-m)!_C\right).$$ 
It is convenient to write $ m_1 \colonequal m=\alpha(p) m_1' + r_1 $, and $m_2\colonequal(n-m)=\alpha(p)m_2'+r_2$, where $ 0 \leq r_1,r_2 < \alpha(p) $ and $ m_1',m_2' \in \N $, and consider
\begin{align*}
\nu_p\!\left(\binom{\alpha(p)n'+r}{\alpha(p)m_1'+r_1}_C\right) &= \nu_p\!\left((\alpha(p)n' + r)!_C\right)-\nu_p\!\left((\alpha(p)m_1' + r_1)!_C\right)-\nu_p\!\left((\alpha(p)m_2' + r_2)!_C\right).
\end{align*}
After substituting Lemma~\ref{lemma:idealCorial}, this becomes
\begin{align*}
    & s(p)\left(\left\lfloor\frac{\alpha(p)n'+r}{\alpha(p)}\right\rfloor-\left\lfloor\frac{\alpha(p)m_1'+r_1}{\alpha(p)}\right\rfloor-\left\lfloor\frac{\alpha(p)m_2' + r_2}{\alpha(p)}\right\rfloor\right) \\ & +\nu_p\left(\left\lfloor\frac{\alpha(p)n'+r}{\alpha(p)}\right\rfloor!\right) - \nu_p\left(\left\lfloor\frac{\alpha(p)m_1'+r_1}{\alpha(p)}\right\rfloor!\right) - \nu_p\left(\left\lfloor\frac{\alpha(p)m_2'+r_2}{\alpha(p)}\right\rfloor!\right).
\end{align*}
Simplifying each floor, this valuation equals
\begin{equation} \label{eq:valuationCnomial}
\nu_p\!\left(\binom{\alpha(p)n'+r}{\alpha(p)m_1'+r_1}_C\right)= s(p)\left(n'-m_1'-m_2'\right) + \nu_p\left(n'!\right) - \nu_p\left(m_1'!\right) - \nu_p\left(m_2'!\right).
\end{equation}
Our analysis now splits based the value of $\lambda:= n'-m_1'-m_2'$. 
Because $n=m_1+m_2$, we may state
$$r+\lambda\alpha(p)=r_1+r_2.$$
Therefore, by the bounds on $r$, $r_1$, and $r_2$, $\lambda$ may only take on the value of $0$ or $1$. 

The first case we consider is $\lambda=0$. In this case, $n'-m_1'-m_2'=0$. Thus Equation ~\eqref{eq:valuationCnomial} simplifies to
\begin{equation*} 
\nu_p\!\left(\binom{\alpha(p)n'+r}{\alpha(p)m_1'+r_1}_C\right)=
\nu_p\left(n'!\right) - \nu_p\left(m_1'!\right) - \nu_p\left(m_2'!\right) = \nu_p\left(\frac{n'!}{m_1'!m_2'!}\right) = \nu_p\!\left(\binom{n'!}{m_1'!}\right).
\end{equation*}
We substitute this for the right hand side of the $p$-adic valuation of the $C$-nomial, finding \begin{align*}
    \nu_p\!\left(\binom{\alpha(p)n'+r}{\alpha(p)m_1'+r_1}_C\right) &= \nu_p\!\left(\binom{n'}{m_1'}\right).
\end{align*}

The other case is $\lambda=1$.
In this case, $n'-m_1'-m_2'=1$. Thus the first addend on the right hand side of Equation ~\eqref{eq:valuationCnomial} evaluates to $s(p)$. Since $m_1'+m_2'=n-1$, notice that the latter addends simplify as 
$$\nu_p\left(n'!\right) - \nu_p\left(m_1'!\right) - \nu_p\left(m_2'!\right)= \nu_p\left(\frac{n' (n'-1)!}{m_1'! m_2'!}\right)= \nu_p\!\left(n'\binom {n'-1} {m_1'}\right).$$
Therefore, the $p$-adic valuation of the $C$-nomial in this case is
\begin{align*}
    \nu_p\!\left(\binom{\alpha(p)n'+r}{\alpha(p)m_1'+r_1}_C\right) &= s(p)+\nu_p\!\left(n'\binom{n'-1}{m_1'}\right).
\end{align*}

Now let's interpret what these cases mean and how they allow us to simplify the summation 
$$\sum_{m=0}^{n} x^{\nu_p\!\left( \binom{n}{m}_C \right)} = \sum_{\alpha(p)m_1'+r_1=0}^{\alpha(p)n'+r} x^{\nu_p\!\left(\binom{\alpha(p)n'+r}{\alpha(p)m+r_1}_C\right)}. $$ 
When $\lambda=0$, $r=r_1+r_2$. Therefore, by the bounds on $r_i$, $r_1\in\{0,1,\ldots,r\}$ and from the iterating variable $\alpha(p)m_1' + r_1$ running up to $\alpha(p)n'+r$, the valid values of $m_1'$ are $0 \leq m_1' \leq n'$. Therefore, this evaluates the summation in this case to $$(r+1)\sum_{m_1'=0}^{n'} x^{\nu_p\!\left(\binom{n'}{m_1'}\right)}.$$
When $\lambda=1$, $r+\alpha(p)=r_1+r_2$. In this case, valid values are $r_1\in\{r+1, \ldots,\alpha(p)-1\}$, and from the iterating variable $\alpha(p)m_1' + r_1$ running up to $\alpha(p)n'+r$, the valid values of $m_1'$ are $0 \leq m_1' \leq n'-1$. Therefore, this evaluates the summation in this case to $$(\alpha(p)-r-1)\sum_{m_1'=0}^{n'-1} x^{s(p)+\nu_p\!\left(n'\binom{n'}{m_1'}\right)}.$$
Adding these two components completes the proof.
\end{proof}

\end{document}
